\chapter*{Resumen}
\addcontentsline{toc}{chapter}{Resumen}

Este documento presenta el desarrollo de un sistema de locomoción para el
robot cuadrúpedo Nova Spot Micro. Se implementaron modelos cinemáticos y
dinámicos nominales, una descripción reproducible para Gazebo y MuJoCo, un
controlador de postura y marchas cartesianas tipo paso y gateo, métricas en
línea y un supervisor de seguridad. La arquitectura se integró en ROS~2 Humble
y se verificó mediante pruebas automatizadas y ensayos registrados con
\textit{rosbag2}.

La línea base de gateo con temporización corregida reprodujo ciclos de
4,32~s con diferencias menores al 0,2~\% entre dos ejecuciones. La marcha paso
completó ensayos independientes en Gazebo y una validación articular en MuJoCo.
La instrumentación simulada incluyó contacto en los cuatro pies, IMU y margen
de estabilidad nominal. Un ensayo específico de 24 ciclos distinguió el
contacto crudo del contacto filtrado: las interrupciones de las patas traseras
duraron menos de 0,09~s y ninguna superó el criterio de persistencia de
0,12~s. Por tanto, el avance observado no demostró un vuelo trasero sostenido y
se concluyó que el patrón físico simulado todavía no coincide con la secuencia
ideal de contactos.

Los resultados descritos corresponden al estado verificable del proyecto. La
caracterización física completa, la seguridad eléctrica, la calibración de los
doce servomotores y la comparación final con una política de aprendizaje por
refuerzo permanecen en desarrollo y no se presentan como logros terminados.

\textbf{Palabras clave:} robot cuadrúpedo, locomoción, ROS~2, Gazebo, MuJoCo,
contacto, estabilidad, aprendizaje por refuerzo.

\chapter*{Abstract}
\addcontentsline{toc}{chapter}{Abstract}

This document presents the development of a locomotion system for the Nova
Spot Micro quadruped robot. Nominal kinematic and dynamic models, reproducible
Gazebo and MuJoCo descriptions, posture and Cartesian gait controllers, online
metrics, and a safety supervisor were implemented using ROS~2 Humble. Baseline
experiments demonstrated repeatable crawl and step trajectories in simulation.
A 24-cycle experiment separated raw from debounced foot-contact observations.
Rear-foot interruptions remained below 0.09~s and never met the 0.12~s
persistence criterion; therefore, simulated forward motion did not demonstrate
sustained rear-foot flight. Physical characterization, electrical safety,
twelve-servo calibration, and the final reinforcement-learning comparison are
still in progress and are not reported as completed results.

\textbf{Keywords:} quadruped robot, locomotion, ROS~2, Gazebo, MuJoCo,
contact, stability, reinforcement learning.
